\chapter{Occupancy Prediction}
\label{chap:occupancy-prediction}

\section{Purpose of Occupancy Prediction}

Occupancy prediction is the first decision-making stage of the smart hotel
application. Its purpose is to estimate the future operating state of a room
before a lighting or HVAC recommendation is produced. In the implemented
system, the output belongs to one of three classes:

\begin{itemize}
    \item \textit{Occupied}: a guest is expected to be present in the room;
    \item \textit{Vacant}: the room is expected to be unoccupied; and
    \item \textit{Cleaning}: housekeeping activity is expected.
\end{itemize}

This distinction is important because the appropriate control strategy is
different for each state. An occupied room should prioritize comfort, a vacant
room can use more aggressive energy-saving settings, and a room under cleaning
may require bright lighting and a moderate HVAC setting for hotel staff.

\section{Data Sources and Prediction Target}

The occupancy subsystem combines records from
\texttt{PIRSensorData.csv} and \texttt{hotelReservationData.csv}. The PIR
dataset supplies timestamped room activity and room-state information. The
reservation dataset provides guest and booking context. The two datasets are
joined using the guest identifier when a valid identifier is available.

The classical occupancy model predicts the room state 60 minutes in the
future. Because measurements are recorded every five minutes, this corresponds
to shifting the target by 12 samples:

\[
    y_t = \mathrm{room\_state}_{t+12}.
\]

This target is intentionally different from simply reproducing the current
room state. It allows the recommendation system to prepare for a state change,
for example by reducing energy use before a room becomes vacant.

\section{Feature Engineering}

The Random Forest training script constructs both instantaneous and historical
features. The numerical features are:

\begin{itemize}
    \item room number;
    \item current PIR motion value;
    \item mean PIR motion during the previous hour;
    \item total number of motion detections during the previous hour;
    \item number of adults and children;
    \item total guest count;
    \item whether an active guest identifier is present;
    \item hour-of-day and day-of-week cyclic features;
    \item total nights, total reservation amount and stay duration; and
    \item price per night.
\end{itemize}

Nationality and room type are treated as categorical variables. Missing
numerical values are replaced with the median, categorical values are replaced
with the most frequent value, and categorical variables are one-hot encoded.
Numerical values are standardized before they are passed to the classifier.

Time is encoded cyclically because hour 23 and hour 0 are temporally close even
though their raw numerical values are far apart:

\[
    h_{\sin} = \sin\left(\frac{2\pi h}{24}\right), \qquad
    h_{\cos} = \cos\left(\frac{2\pi h}{24}\right),
\]

\[
    d_{\sin} = \sin\left(\frac{2\pi d}{7}\right), \qquad
    d_{\cos} = \cos\left(\frac{2\pi d}{7}\right).
\]

\section{Random Forest Occupancy Model}

The classical model is implemented as a scikit-learn pipeline containing a
\texttt{ColumnTransformer} and a \texttt{RandomForestClassifier}. The forest
uses 160 trees, a minimum leaf size of 20, balanced subsample class weights,
parallel execution and a fixed random seed of 42. Class weighting is necessary
because Cleaning observations are much less frequent than Occupied and Vacant
observations.

A chronological 80/20 split is used instead of a random split. Earlier
observations are used for training and later observations are used for testing.
This prevents information from future timestamps from leaking into the
training set and gives a more realistic estimate of deployment performance.

\section{Transformer Occupancy Model}

The Transformer implementation uses sequences of 12 five-minute observations,
representing one hour of room history, and predicts the state 12 steps later.
Each time step contains 17 features:

\begin{itemize}
    \item PIR motion, adult count, child count and total guest count;
    \item active-guest status;
    \item cyclic hour and day-of-week values;
    \item normalized reservation duration and cost variables; and
    \item one-hot room-type indicators for Deluxe, Standard, Suite and Unknown.
\end{itemize}

The model first projects the 17-dimensional input to a 64-dimensional hidden
space. A trainable positional embedding is added so the network can distinguish
the order of observations. The encoder contains two Transformer layers with
four attention heads, a feed-forward dimension of 128, GELU activation and
dropout of 0.15. After encoding, the hidden states are normalized and averaged
over time. A fully connected classification head produces logits for Cleaning,
Occupied and Vacant.

The model is trained with AdamW and weighted cross-entropy loss. Gradient
clipping is applied during training to reduce instability. The training set
contains 320,000 sequences and the test set contains 80,000 sequences.

\section{Serving-Time Prediction Workflow}

The Flask endpoint \texttt{/api/predict\_occupancy} accepts a room number,
timestamp, model type and optional lookback and prediction-horizon values. The
backend queries the \texttt{pir\_sensor\_data} table for records between the
lookback start time and requested timestamp. The user interface allows the
Random Forest, Transformer or automatic model-selection mode to be chosen.

For a Transformer request, the recent observations are converted into a
fixed-length sequence and passed through the saved PyTorch model. A softmax
operation converts logits into probabilities. For a Random Forest request, a
single engineered feature row is passed to the saved scikit-learn pipeline.
Both paths return:

\begin{itemize}
    \item the predicted class;
    \item prediction confidence;
    \item a probability for each occupancy class;
    \item diagnostic features such as motion rate and latest room state; and
    \item the model name and prediction horizon.
\end{itemize}

The backend includes two fallback strategies. If no recent PIR history exists,
it returns a conservative default distribution that favors Vacant. If a saved
model is unavailable, a rule-based score combines the latest state, active
guest status, guest count, motion rate and motion count. These fallbacks keep
the API operational, although model-based predictions are preferred.

\section{Occupancy Results}

The Random Forest was evaluated on 1,042,300 test rows and achieved an overall
accuracy of 98.9\%. The Transformer achieved 98.8\% accuracy on 80,000 test
sequences. Occupied and Vacant were predicted with high F1-scores by both
models. However, Cleaning remained difficult because of its very small support.
The Random Forest obtained a Cleaning F1-score of 0.100, while the Transformer
improved it to 0.248.

These results show why overall accuracy must be interpreted together with
class-level precision, recall and F1-score. A model can obtain high accuracy
on an imbalanced dataset while still performing poorly on an operationally
important minority class.


\chapter{Guest Persona Classification}
\label{chap:persona-classification}

\section{Role of Persona Classification}

Persona classification represents repeated guest behavior as an interpretable
category. The application implements separate persona models for lighting and
temperature because brightness preference and thermal preference are not
necessarily the same. Persona outputs are not final control commands; they are
context variables used by the recommendation models.

\section{Lighting Persona Categories}

The lighting subsystem uses the following classes:

\begin{itemize}
    \item \textit{Balanced}: moderate and varied lighting use;
    \item \textit{Routine}: repeated lighting patterns at similar times;
    \item \textit{StaticBright}: consistently high preferred brightness;
    \item \textit{StaticDim}: consistently low preferred brightness;
    \item \textit{NightFocused}: stronger activity during late-night hours;
    \item \textit{Housekeeping}: bright task-oriented lighting; and
    \item \textit{Unknown}: insufficient or unsupported history.
\end{itemize}

The five primary guest personas are learned by the trained classifiers.
Housekeeping and Unknown are also recognized by the live recommendation
service so that operational and missing-data cases can be handled safely.

\section{Random Forest Lighting Persona Model}

The classical lighting persona script groups lighting records by room and
date, producing one feature row for each room-day. Only records with an active
reservation are used. The target is the dominant
\texttt{lightning\_persona} value for that room-day.

The aggregated features include:

\begin{itemize}
    \item mean, standard deviation and maximum lighting level;
    \item percentage of records in which a lamp is on;
    \item PIR motion rate;
    \item mean and maximum occupant count;
    \item mean number of active actors;
    \item average lighting value for each hour of the day; and
    \item the proportional usage of each lamp location.
\end{itemize}

Missing values are imputed with the median and the numerical features are
standardized. The Random Forest contains 200 trees, uses a minimum leaf size of
five and applies balanced subsample class weighting. A stratified 75/25 split
preserves class proportions when each class has sufficient examples.

\section{Transformer Lighting Persona Model}

The Transformer treats each room-day as a complete sequence of 288 five-minute
time slots. At every slot, the model receives 19 features: the normalized level
of each of the 11 lamp locations, PIR motion, number of occupants, active
actors, three behavioral indicators, and sine/cosine hour encodings.

Missing time slots are reindexed and filled with zero. This creates a regular
daily sequence even if a lamp has no recorded activity during some periods.
Brightness values are clipped to the operational range of 0--80 and normalized
by dividing by 80.

The model projects the input into a 96-dimensional representation. It uses
three Transformer encoder layers, four attention heads, a feed-forward
dimension of 192, GELU activation and dropout of 0.15. After layer
normalization, mean pooling summarizes the full day. A classification head
maps the pooled representation to persona logits.

Training uses AdamW, class-weighted cross-entropy and gradient clipping. The
saved model was trained on 9,132 room-day samples, of which 7,305 were used for
training and 1,827 for testing.

\section{Live Lighting Persona Prediction}

The endpoint \texttt{/api/predict\_lighting\_persona} retrieves recent records
from the \texttt{lightning\_data} table. The default history window is 24
hours, while the user interface also supports shorter three-hour and six-hour
contexts.

Before prediction, the backend calculates diagnostics that are displayed in
the interface:

\begin{itemize}
    \item mean level while lamps are on;
    \item proportion of lamp records with a level greater than zero;
    \item mean nighttime and daytime levels;
    \item nighttime share of total day/night brightness; and
    \item mean number of active actors.
\end{itemize}

The Random Forest path constructs the same room-day style feature row used
during training. The Transformer path builds a 288-step sequence aligned to the
requested timestamp. Both return class probabilities and confidence. If no
trained model is available, a rule-based fallback scores personas using recent
mean brightness, lighting ratio, night share and active-actor information.

\section{Temperature Persona Categories}

Temperature behavior is classified into:

\begin{itemize}
    \item \textit{AlwaysOnComfort}: prioritizes a narrow comfort range;
    \item \textit{EnergySaver}: accepts wider temperature variation;
    \item \textit{Reactive}: changes settings after conditions become
    uncomfortable;
    \item \textit{Preconditioning}: prepares the room before occupancy; and
    \item \textit{Housekeeping}: represents staff-oriented HVAC operation.
\end{itemize}

\section{Temperature Persona Transformer}

The temperature persona model uses 24 sequential observations. Because the
temperature dataset is sampled every five minutes, the sequence represents two
hours of recent room behavior. The input contains 42 encoded features,
including:

\begin{itemize}
    \item scaled floor, room size, outside temperature, room temperature,
    setpoint and ideal temperature;
    \item temperature error and comfort error;
    \item PIR motion and guest-presence indicators;
    \item cyclic hour and day-of-week features;
    \item room facade and room type;
    \item HVAC mode, occupant state and PIR persona; and
    \item current room state.
\end{itemize}

The model uses a 96-dimensional hidden representation, four attention heads,
two encoder layers, a feed-forward dimension of 192 and dropout of 0.15.
Training uses class-weighted cross-entropy, AdamW and a chronological
train/validation/test split. From 500,000 generated sequences, 400,000 were
used for training, 50,000 for validation and 50,000 for testing.

\section{Persona Classification Results}

The Random Forest lighting persona classifier achieved 88.5\% accuracy and a
weighted F1-score of 0.885. The Transformer improved the lighting persona
accuracy to 92.5\%, with a macro F1-score of 0.926. StaticDim achieved the
highest Transformer F1-score, 0.952.

The temperature persona Transformer achieved 80.55\% accuracy and a weighted
F1-score of 0.8088. AlwaysOnComfort achieved an F1-score of 0.9006, while
EnergySaver and Reactive were more frequently confused. Housekeeping and
Preconditioning had very high scores but very small support, so those results
should not be interpreted as evidence of broad generalization.


\chapter{Lighting Recommendation System}
\label{chap:lighting-recommendation}

\section{Recommendation Pipeline}

The lighting recommendation system is implemented as a chained pipeline:

\begin{enumerate}
    \item obtain the room number, timestamp and history window;
    \item predict future occupancy;
    \item predict the guest's lighting persona;
    \item retrieve recent room and per-lamp history;
    \item generate one recommendation for each known lamp;
    \item enforce hardware and operational constraints; and
    \item convert current and recommended levels into energy estimates.
\end{enumerate}

The Flask endpoint \texttt{/api/lightning\_recommendation} can receive
occupancy and persona values directly, or it can generate them automatically
when \texttt{use\_model\_predictions} is enabled. In automatic mode, the
occupancy model uses a one-hour lookback and a 60-minute horizon. The persona
model uses the selected lighting history window.

\section{Lamp Classification}

The implementation distinguishes between dimmable LEDs and non-dimmable bulbs.
The dimmable group contains:

\begin{itemize}
    \item \texttt{bed\_left}, \texttt{bed\_right},
    \texttt{dinner\_table}, \texttt{table} and \texttt{hidden\_top}.
\end{itemize}

The non-dimmable group contains:

\begin{itemize}
    \item \texttt{closet}, \texttt{corridor\_left},
    \texttt{corridor\_right}, \texttt{shower}, \texttt{cabinet} and
    \texttt{sink}.
\end{itemize}

This classification affects both recommendation and energy calculation.
Dimmable lamps may receive a continuous level between 0 and 80. Non-dimmable
lamps are constrained to either 0 or 80.

\section{Historical Context}

For each lamp, the backend obtains the latest level and the mean values over
the previous one hour, three hours and full lookback period. Room-wide
context includes floor, PIR motion, number of occupants, active actors,
reservation status and behavioral indicators such as
\texttt{hurry\_morning}, \texttt{lazy\_day} and \texttt{forgetful}.

The tabular recommendation model additionally uses rolling features:

\begin{itemize}
    \item per-lamp mean level over 1, 3 and 24 hours;
    \item per-lamp on-rate over three hours;
    \item room-wide mean level over 1 and 24 hours; and
    \item room-wide motion rate over one hour.
\end{itemize}

All rolling features are shifted by one sample during training. This ensures
that a training target is predicted only from earlier observations and avoids
using the current target value as historical input.

\section{Efficient Recommendation Target}

The recommendation target is not a direct copy of historical brightness.
Instead, the training script constructs an energy-aware target using occupancy,
reservation state, recent preference and persona policy.

If the room is Vacant or has no active reservation, the target is zero. If the
room is being cleaned, an active lamp is assigned level 80. If a lamp was
already off, it remains off. Otherwise, the target is based on the minimum of
the current level, recent median preference and persona cap, multiplied by a
persona reduction factor:

\[
    L_{\mathrm{raw}} =
    \min(L_{\mathrm{current}},L_{\mathrm{recent}},C_p)F_p.
\]

The result is then limited by the persona's minimum active level:

\[
    L_{\mathrm{rec}} =
    \mathrm{clip}\left(
    \max(L_{\mathrm{raw}},M_p),0,80\right).
\]

For example, StaticBright uses a factor of 0.90, a cap of 75 and a minimum
active level of 35. StaticDim uses a factor of 0.62, a cap of 35 and a minimum
of 10. NightFocused uses a factor of 0.68 and a cap of 45. These policies lower
avoidable brightness while retaining differences between guest preferences.

\section{HistGradientBoosting Recommendation Model}

The tabular model uses a \texttt{HistGradientBoostingRegressor} with 180
iterations, a learning rate of 0.08, 31 maximum leaf nodes and L2
regularization of 0.01. Numerical variables are median-imputed and standardized.
Categorical variables---lamp location, reservation status, occupancy prediction
and persona prediction---are imputed and one-hot encoded.

Training uses 10,563,249 rows. The first 80\% of the timeline is used for
training and the final 20\% for testing. Predictions are rounded and clipped
to the 0--80 operational range.

\section{Transformer Recommendation Model}

The Transformer regressor uses a sequence of 12 recent five-minute records and
predicts the efficient level for the next five-minute step. Each record
contains 38 encoded features. These include normalized brightness, PIR motion,
occupant counts, behavioral flags, cyclic time values, lamp location,
reservation status, occupancy prediction and lighting persona prediction.

The model has a 96-dimensional hidden space, four attention heads, two encoder
layers, a feed-forward dimension of 192 and dropout of 0.15. A sinusoidal
positional encoding is added to the projected input. The representation of the
last sequence token is passed through a regression head. A sigmoid constrains
the normalized prediction before it is converted back to the 0--80 level
range.

\section{Serving Constraints and Fallback Behavior}

The live backend can select the HistGradientBoosting model, Transformer model
or automatic mode. If the selected model files are unavailable, the service
uses the explicit persona policy as a rule-based fallback.

Several constraints are applied after model inference:

\begin{itemize}
    \item a currently inactive lamp remains off;
    \item a non-dimmable lamp is snapped to either 0 or 80;
    \item a Vacant room or inactive reservation turns lamps off;
    \item Cleaning mode keeps required lamps bright; and
    \item recommendations remain within the valid device range.
\end{itemize}

These constraints are important because a statistically valid model output
may still be physically impossible or operationally inappropriate.

\section{Recommendation Output}

The API returns a record for every lamp containing current level, recommended
level, one-, three- and 24-hour means, lamp type, estimated current and
recommended energy, saved energy, last observation time, model name and a
human-readable explanation.

The React interface displays current and recommended bars for each lamp,
prediction confidence, the selected model, total dimmable energy, estimated
energy after recommendation, saving relative to current use and saving relative
to a full-brightness baseline.


\chapter{Temperature and HVAC Recommendation}
\label{chap:temperature-recommendation}

\section{Overview}

The temperature subsystem first predicts a temperature persona and then
generates an HVAC setpoint. The endpoint
\texttt{/api/tempreture\_recomendation} follows this chain automatically when
model predictions are requested. The spelling in the API and filenames is
\texttt{tempreture\_recomendation}; however, the function represents a
temperature recommendation.

\section{Input Sequence}

The recommendation service queries the \texttt{temperature\_data} table for
the selected room and lookback interval. The default history window is two
hours. The sequence includes floor, facade, room type, room size, outside
temperature, indoor temperature, current setpoint, ideal temperature, HVAC
mode, occupant state, PIR persona, room state, PIR motion and guest presence.

The final encoded input has 49 features because categorical variables are
expanded into one-hot indicators. Two upstream prediction fields are included:

\begin{itemize}
    \item occupancy prediction: Cleaning, Occupied, Vacant or Unknown; and
    \item temperature persona prediction: AlwaysOnComfort, EnergySaver,
    Housekeeping, Preconditioning, Reactive or Unknown.
\end{itemize}

\section{Construction of the Training Target}

The target setpoint is generated from an explicit comfort and efficiency
policy. First, an operating mode is inferred. An existing heating or cooling
mode is retained. Otherwise, cooling is selected when the room is at least
$0.7\,^{\circ}\mathrm{C}$ above the ideal temperature or the outside
temperature is at least $27\,^{\circ}\mathrm{C}$. Heating is selected when the
room is at least $0.7\,^{\circ}\mathrm{C}$ below the ideal temperature or the
outside temperature is at most $12\,^{\circ}\mathrm{C}$. All remaining cases
use idle mode.

Persona policies define nominal setpoints. For example:

\begin{itemize}
    \item AlwaysOnComfort: $22.5\,^{\circ}\mathrm{C}$ cooling and
    $22.0\,^{\circ}\mathrm{C}$ heating;
    \item EnergySaver: $26.0\,^{\circ}\mathrm{C}$ cooling and
    $19.0\,^{\circ}\mathrm{C}$ heating;
    \item Reactive: $23.5\,^{\circ}\mathrm{C}$ cooling and
    $21.5\,^{\circ}\mathrm{C}$ heating; and
    \item Housekeeping: $24.0\,^{\circ}\mathrm{C}$ cooling and
    $20.0\,^{\circ}\mathrm{C}$ heating.
\end{itemize}

Vacant rooms use more energy-saving targets: $28\,^{\circ}\mathrm{C}$ for
cooling, $17\,^{\circ}\mathrm{C}$ for heating and
$25\,^{\circ}\mathrm{C}$ for idle operation.

To prevent an unrealistic change in a single five-minute step, the target is
smoothed using the policy target, ideal temperature and current setpoint:

\[
    T_{\mathrm{target}} =
    0.65T_{\mathrm{policy}} +
    0.25T_{\mathrm{ideal}} +
    0.10T_{\mathrm{current}}.
\]

The result is rounded to the nearest $0.5\,^{\circ}\mathrm{C}$ and clipped to
the interval from $16\,^{\circ}\mathrm{C}$ to $28\,^{\circ}\mathrm{C}$.

\section{Transformer Regressor}

The model processes 24 observations, corresponding to two hours of five-minute
samples. It uses a 96-dimensional projection, sinusoidal positional encoding,
two Transformer encoder layers, four attention heads, a feed-forward dimension
of 192, GELU activation and dropout of 0.15.

The last encoded token is passed through layer normalization and a two-layer
regression head. A sigmoid limits the normalized output to the interval
$[0,1]$. It is then transformed back into degrees Celsius:

\[
    \hat{T} = 16 + 12\hat{y}.
\]

Training uses Smooth L1 loss, AdamW, a learning-rate scheduler and chronological
training, validation and test partitions. A total of 500,000 sequences were
used: 400,000 for training, 50,000 for validation and 50,000 for testing.

\section{Live Recommendation and Explanation}

At serving time, the predicted setpoint is clipped and rounded to the nearest
half degree. The backend compares it with the current setpoint and labels the
action as \textit{keep}, \textit{raise} or \textit{lower}. It also infers a
target mode. In the live implementation, the current HVAC mode is retained
when it is heating or cooling; otherwise, room-to-setpoint and outdoor
temperature thresholds determine heating, cooling or idle operation.

The API response includes the current setpoint, recommended setpoint, delta,
target mode, action, model name, input predictions and a textual reason. The
reason explicitly notes when Vacant status permits a more efficient setpoint,
when EnergySaver allows a wider comfort band, or when AlwaysOnComfort gives
greater priority to comfort.

\section{Immediate Energy Estimate}

The recommendation service estimates the power required by both the current
and recommended setpoints using the same HVAC thermal-demand model used by the
energy report. For the latest five-minute sample:

\[
    E_{\mathrm{current}} =
    P_{\mathrm{current}}\frac{5}{60},
    \qquad
    E_{\mathrm{recommended}} =
    P_{\mathrm{recommended}}\frac{5}{60}.
\]

The API returns current and recommended power, current and recommended
watt-hours, absolute saving and percentage saving. These values are model-based
estimates rather than direct meter readings.


\chapter{Energy Consumption Analysis}
\label{chap:energy-analysis}

\section{Lighting Power Model}

Lighting energy is calculated from
\texttt{EnergyConsumption/DimmerGraphs.xlsx}. The application reads the
numeric measurement block and extracts dimmer level, lux, voltage and current
for LEDs and bulbs.

LED current is recorded in microamperes, while bulb current is recorded in
milliamperes. The implementation therefore uses:

\[
    P_{\mathrm{LED}} =
    \frac{V_{\mathrm{LED}}I_{\mathrm{LED},\mu A}}{10^6},
\]

\[
    P_{\mathrm{bulb}} =
    \frac{V_{\mathrm{bulb}}I_{\mathrm{bulb},mA}}{10^3}.
\]

The application records levels on a 0--100 scale, but the physical dimmer table
uses a 0--254 scale. A recorded application level is converted using:

\[
    L_{\mathrm{device}} =
    \mathrm{round}\left(L_{\mathrm{app}}\frac{254}{100}\right).
\]

The maximum operational application level is 80. Dimmable LEDs use the power
looked up for their converted level. Non-dimmable bulbs use full power whenever
their recorded level is greater than zero.

\section{Lighting Energy Aggregation}

The endpoint \texttt{/api/energy} receives a room and date interval. The
database query selects active lamp records, excludes the artificial
\texttt{none} location and includes the start timestamp while excluding the
end timestamp.

For each five-minute record:

\[
    E_i = P_i\frac{5}{60}.
\]

The system calculates actual energy and a maximum-level baseline for every
record. It then aggregates:

\begin{itemize}
    \item total actual and maximum energy;
    \item absolute and percentage saving;
    \item energy grouped by lamp;
    \item a dimmable-only summary; and
    \item the number of active samples.
\end{itemize}

Savings from dimming are calculated separately because non-dimmable bulbs have
equal actual and maximum power whenever they are on. The API also evaluates
the AI lighting recommendation over every timestamp in the selected interval.
This comparison reports recommended energy against a full-brightness dimmable
baseline.

\section{HVAC Thermal-Demand Model}

Direct HVAC electricity measurements are not present in the dataset.
Consequently, the system estimates HVAC power from thermal demand. The rated
limits are 2000 W for heating and 1500 W for cooling.

For heating:

\[
    \Delta T_{\mathrm{room}} =
    \max(T_{\mathrm{set}}-T_{\mathrm{room}},0),
\]

\[
    \Delta T_{\mathrm{weather}} =
    \max(T_{\mathrm{set}}-T_{\mathrm{outside}},0),
\]

\[
    P_{\mathrm{raw}} =
    420 + 260\Delta T_{\mathrm{room}}
    + 32\Delta T_{\mathrm{weather}}.
\]

For cooling:

\[
    \Delta T_{\mathrm{room}} =
    \max(T_{\mathrm{room}}-T_{\mathrm{set}},0),
\]

\[
    \Delta T_{\mathrm{weather}} =
    \max(T_{\mathrm{outside}}-T_{\mathrm{set}},0),
\]

\[
    P_{\mathrm{raw}} =
    380 + 230\Delta T_{\mathrm{room}}
    + 30\Delta T_{\mathrm{weather}}.
\]

Room size is normalized against a reference size of $35\,\mathrm{m}^2$ and
clipped between 0.70 and 1.45. Occupancy factors are 0.55 for Vacant, 0.75 for
Cleaning, 0.95 for Occupied without motion, and 1.10 for Occupied with motion.
Unknown states use 0.85 without motion and 1.00 with motion.

The final estimate is:

\[
    P_{\mathrm{HVAC}} =
    \mathrm{clip}
    \left(P_{\mathrm{raw}}F_{\mathrm{size}}F_{\mathrm{occupancy}},
    0,P_{\mathrm{rated}}\right).
\]

Off and unsupported modes produce zero power. Every sample contributes
$P_{\mathrm{HVAC}}(5/60)$ watt-hours.

\section{HVAC Baseline and Reports}

The HVAC baseline represents continuous operation at rated power for the
dominant active mode. The dominant mode is whichever of heating or cooling
appears in more samples. If neither appears, heating is used as a fallback.
For $N$ samples:

\[
    E_{\mathrm{max}} =
    P_{\mathrm{rated,dominant}}N\frac{5}{60}.
\]

The endpoint \texttt{/api/hvac\_energy} returns total, heating and cooling
energy, active and total minutes, mode-level averages, maximum baseline,
estimated savings and AI-recommendation energy. The frontend presents actual
versus baseline energy, energy by mode, time-share charts and an hourly or
daily energy time series.

\section{Interpretation}

Lighting energy uses measured voltage-current relationships and is therefore
closer to a physical calculation. HVAC energy is a model-based estimate and
should be interpreted as comparative rather than meter-grade consumption.
Both analyses assume that one database row represents exactly five minutes.
Missing samples or irregular sensor intervals can therefore affect totals.


\chapter{Machine Learning Models}
\label{chap:machine-learning-models}

\section{Model Families}

The application uses three principal model families:

\begin{itemize}
    \item Random Forest classifiers for tabular occupancy and lighting persona
    prediction;
    \item HistGradientBoosting regression for tabular lighting recommendation;
    and
    \item Transformer encoders for sequential occupancy, persona, lighting and
    temperature tasks.
\end{itemize}

The classical models provide efficient and interpretable baselines. The
Transformers are used when the order and temporal interaction of observations
are expected to contain useful information.

\section{Preprocessing Pipelines}

The scikit-learn models save preprocessing and prediction as a single pipeline.
Numerical variables are imputed and standardized. Categorical variables are
imputed and one-hot encoded with unknown-category handling. Saving the full
pipeline ensures that training-time transformations are repeated consistently
during API inference.

Transformer inputs use explicit scaling and one-hot feature construction saved
in metadata files. The serving code reads the metadata to reconstruct feature
order, class labels, sequence length and model dimensions. This is important
because changing the order of encoded columns would make a trained neural
network invalid even if the same values were supplied.

\section{Transformer Architecture}

All Transformer models follow a related encoder-only design:

\begin{enumerate}
    \item project each input vector into a hidden representation;
    \item add trainable or sinusoidal positional information;
    \item process the sequence using multi-head self-attention and feed-forward
    layers;
    \item normalize and pool the sequence or select the final token; and
    \item apply a classification or regression head.
\end{enumerate}

Classification models use softmax probabilities and weighted cross-entropy.
Regression models produce a normalized scalar that is converted to either a
lighting level or HVAC setpoint. AdamW is used throughout the Transformer
training scripts. Dropout and weight decay reduce overfitting, while gradient
clipping is used in the classification scripts.

\section{Class Imbalance}

Class imbalance is addressed through balanced class weights in Random Forest
models and weighted cross-entropy in Transformer classifiers. The class weight
for class $c$ is calculated from its training frequency:

\[
    w_c = \frac{N}{K N_c},
\]

where $N$ is the number of training examples, $K$ is the number of classes and
$N_c$ is the number of examples belonging to class $c$.

This increases the penalty for misclassifying rare classes. Nevertheless, the
Cleaning occupancy results show that weighting alone cannot fully compensate
for very limited examples.

\section{Training and Data Splitting}

Most time-dependent models use chronological splitting. Sequence end times are
sorted, and earlier sequences are assigned to training while later sequences
are assigned to validation and testing. This is more appropriate than a random
split for time-series applications because it simulates prediction on future
data.

The lighting persona Random Forest uses a stratified random split because its
examples are aggregated room-days rather than direct next-step forecasts. All
training scripts use seed 42 to improve reproducibility.

\section{Evaluation Metrics}

Classification models are evaluated with accuracy, precision, recall,
F1-score and confusion matrices:

\[
    \mathrm{Precision}=\frac{TP}{TP+FP},
    \qquad
    \mathrm{Recall}=\frac{TP}{TP+FN},
\]

\[
    F_1 =
    2\frac{\mathrm{Precision}\cdot\mathrm{Recall}}
    {\mathrm{Precision}+\mathrm{Recall}}.
\]

Regression models use mean absolute error, root mean squared error and
coefficient of determination:

\[
    \mathrm{MAE}=\frac{1}{n}\sum_{i=1}^{n}|y_i-\hat{y}_i|,
\]

\[
    \mathrm{RMSE}=
    \sqrt{\frac{1}{n}\sum_{i=1}^{n}(y_i-\hat{y}_i)^2},
\]

\[
    R^2 =
    1-
    \frac{\sum_i(y_i-\hat{y}_i)^2}
    {\sum_i(y_i-\bar{y})^2}.
\]

The lighting recommendation scripts additionally compare the sum of predicted
levels with actual levels. This is an energy proxy rather than a direct
watt-hour measurement. The application performs the physical watt-hour
calculation separately using the dimmer power table.

\section{Model Persistence and Runtime Selection}

Scikit-learn models are stored with Joblib. Transformer checkpoints are stored
as PyTorch \texttt{.pt} files accompanied by JSON metadata. The backend caches
loaded models and reloads them when the file modification time changes.

The frontend exposes model selectors for occupancy, lighting persona and
lighting recommendation. The backend normalizes aliases such as
\texttt{rf}, \texttt{random\_forest}, \texttt{torch} and
\texttt{transformer}. Automatic mode selects a Transformer when the expected
checkpoint and metadata are available; otherwise it selects the corresponding
classical model.


\chapter{Results and Evaluation}
\label{chap:results}

\section{Summary of Saved Results}

\begin{table}[ht]
\centering
\caption{Evaluation results produced by the project training scripts}
\label{tab:project-model-results}
\begin{tabular}{|p{4.2cm}|p{4.0cm}|p{5.0cm}|}
\hline
\textbf{Task} & \textbf{Model} & \textbf{Saved test result} \\ \hline
Occupancy prediction &
Random Forest &
Accuracy 98.9\%; macro F1 0.692 \\ \hline
Occupancy prediction &
Transformer &
Accuracy 98.8\%; macro F1 0.743 \\ \hline
Lighting persona &
Random Forest &
Accuracy 88.5\%; macro F1 0.893 \\ \hline
Lighting persona &
Transformer &
Accuracy 92.5\%; macro F1 0.926 \\ \hline
Temperature persona &
Transformer &
Accuracy 80.55\%; macro F1 0.8848 \\ \hline
Lighting recommendation &
HistGradientBoosting &
MAE 0.105; RMSE 0.356; $R^2$ 0.999 \\ \hline
Lighting recommendation &
Transformer &
MAE 1.027; RMSE 1.880; $R^2$ 0.985 \\ \hline
Temperature recommendation &
Transformer &
MAE $0.217\,^{\circ}\mathrm{C}$; RMSE
$0.310\,^{\circ}\mathrm{C}$; $R^2$ 0.962 \\ \hline
\end{tabular}
\end{table}

\section{Occupancy Evaluation}

The occupancy models perform very well on Occupied and Vacant observations.
The Random Forest achieved F1-scores of 0.997 and 0.978 for these classes,
respectively. The Transformer achieved 0.991 and 0.990. The Transformer had a
slightly lower overall accuracy but a better macro F1-score because it improved
the Cleaning class.

Cleaning precision remains low: 0.057 for the Random Forest and 0.177 for the
Transformer. This means many Cleaning predictions are false positives.
Cleaning recall is higher, approximately 0.40 for both models, which indicates
that class weighting encourages the models to find Cleaning examples but at
the cost of precision.

\section{Persona Evaluation}

The lighting persona Transformer outperformed the Random Forest in overall
accuracy by four percentage points. Its F1-scores ranged from 0.898 for
Balanced to 0.952 for StaticDim. The improvement suggests that complete
five-minute daily sequences capture temporal patterns that are partly lost
when a day is summarized into aggregated statistics.

The temperature persona model achieved lower overall accuracy. EnergySaver had
high recall of 0.8928 but precision of 0.6568, indicating that the model often
assigned EnergySaver to observations belonging to another class. Reactive had
precision of 0.8291 and recall of 0.7220. These two classes likely overlap
because both can produce setpoint changes in response to room conditions.

\section{Recommendation Evaluation}

The HistGradientBoosting lighting model was tested on 2,112,417 rows. Its
predicted level-sum proxy was 58,723,973 compared with an efficient target of
58,691,230. The target represented a 27.33\% reduction from actual levels, and
the model represented a 27.29\% reduction.

The lighting Transformer used 500,000 sequences, including 50,000 test
sequences. Its model level-sum proxy represented a 29.18\% reduction from
actual levels, compared with 28.96\% for the efficient target. Its errors were
larger than those of the tabular model, although an MAE of approximately one
level remains small relative to the 0--80 range.

The temperature recommendation Transformer achieved an MAE of
$0.217\,^{\circ}\mathrm{C}$. Since the live application rounds the output to
the nearest $0.5\,^{\circ}\mathrm{C}$, this error is smaller than one control
step on average.

\section{End-to-End Evaluation Considerations}

The component metrics evaluate models against labels or policy-derived targets.
They do not by themselves prove real guest comfort or building-level energy
savings. In live operation, recommendation quality also depends on upstream
prediction accuracy. An incorrect occupancy or persona prediction can alter
the final lighting or temperature recommendation.

The application partially addresses this issue by returning complete
probability distributions and diagnostic features to the user interface. It
also permits manual occupancy and persona selection, allowing model-generated
and user-selected scenarios to be compared.


\chapter{Limitations and Future Improvements}
\label{chap:limitations}

\section{Sensor Limitations}

PIR sensors detect movement rather than continuous physical presence. A guest
who is sleeping, reading or otherwise motionless may generate no PIR activity.
Conversely, housekeeping or temporary movement may appear similar to guest
activity. Future work could combine PIR with door sensors, key-card events,
Bluetooth presence, camera-free thermal sensors or CO$_2$ measurements.

\section{Class Imbalance and Label Quality}

Cleaning has far fewer examples than Occupied and Vacant. The existing class
weights improve recall but do not solve low precision. More cleaning examples,
explicit housekeeping schedules, focal loss, balanced sequence sampling or a
two-stage occupancy classifier could improve this class.

Persona labels are derived from dataset behavior and may not correspond to
preferences explicitly stated by real guests. Future studies should collect
guest feedback and measure whether a predicted persona remains stable across
multiple stays.

\section{Policy-Derived Recommendation Targets}

The lighting and temperature recommendation models learn targets generated by
hand-designed efficiency policies. Their high $R^2$ values therefore indicate
that the models reproduce those policies accurately; they do not demonstrate
that the policies are globally optimal.

A future system could learn from explicit comfort feedback, room-service
requests, manual overrides and measured energy cost. Multi-objective
optimization or reinforcement learning could then balance comfort, energy,
equipment wear and electricity price.

\section{Energy Model Limitations}

Lighting energy assumes each record covers exactly five minutes and that the
dimmer lookup table correctly represents installed hardware. Interpolation is
not performed when a converted device level is absent, which can produce
missing power values.

HVAC consumption is estimated rather than measured. The thermal model does not
directly represent humidity, solar gain, wall insulation, door or window
opening, equipment coefficient of performance, fan stages or thermal inertia.
Smart plugs or building-management meters should be used to calibrate and
validate the power model.

\section{Data and Generalization}

The data covers a limited set of room types, personas, weather conditions and
operating periods. Performance may change in another hotel, climate or sensor
installation. Cross-hotel testing and seasonal validation are required before
general deployment.

The Transformer datasets are capped at 400,000 or 500,000 sequences to control
memory use. Although these are large samples, the cap may alter the
distribution if sequence construction reaches some rooms or periods earlier
than others. Stratified room sampling would provide stronger coverage.

\section{Runtime and System Improvements}

Future implementation work should include:

\begin{itemize}
    \item real-time ingestion from an IoT message broker;
    \item authentication and role-based access control;
    \item model and dataset version tracking;
    \item uncertainty thresholds that defer low-confidence decisions;
    \item automatic detection of missing or irregular sensor samples;
    \item explanations based on feature importance or attention analysis;
    \item monitoring for data drift and model degradation;
    \item user override logging and online personalization; and
    \item direct integration with lighting and HVAC controllers.
\end{itemize}

\section{Privacy and Safety}

Reservation and behavioral data can reveal occupancy patterns and guest
preferences. A production system should minimize personally identifiable
information, apply encryption, enforce retention limits and separate identity
from behavioral features. Automatic HVAC control should also use hard safety
limits and permit staff or guests to override recommendations.


\chapter{Conclusion}
\label{chap:conclusion}

This project implements a complete smart hotel analytics and recommendation
pipeline rather than an isolated visualization or prediction model. The React
frontend communicates with a Flask API, which retrieves historical data from
SQLite, executes machine-learning models, applies hardware constraints,
calculates energy and returns explanations and charts.

Occupancy prediction uses PIR and reservation context to estimate whether a
room will be Occupied, Vacant or under Cleaning. Lighting and temperature
persona models convert recent behavior into preference categories. These
outputs are passed to separate recommendation models that generate per-lamp
brightness levels and HVAC setpoints.

The implementation supports both classical tabular models and Transformer
sequence models. Random Forest and HistGradientBoosting models provide strong
baselines, while Transformers capture temporal behavior over one-hour,
two-hour or full-day sequences. The frontend allows model selection and
displays probabilities, diagnostics and model names, making the result more
transparent than a single unexplained command.

Energy analysis is integrated into the recommendation workflow. Lighting
consumption uses measured voltage and current from a dimmer table, with
separate handling for dimmable and non-dimmable devices. HVAC consumption is
estimated using room temperature, setpoint, outside temperature, room size,
occupancy and motion. Both systems compare actual or recommended operation
with an explicit baseline.

The saved experiments show high performance for major occupancy classes,
strong lighting persona classification and low recommendation errors. The
results also reveal important limitations, particularly minority-class
occupancy performance, policy-derived targets and the lack of direct HVAC
metering.

Overall, the project demonstrates how hotel reservation data, IoT sensor
history, machine learning, recommendation policies and energy calculations can
be combined into one decision-support application. The next stage should move
from offline and simulated evaluation toward calibrated physical measurements,
guest feedback, privacy-aware deployment and controlled trials in real hotel
rooms.
